\section*{Chapitre \ref{ch:m1:pnl-and-positions} --- P\&L et comptabilisation d'une position}

\begin{solution}{exo:m1:pnl-and-positions:1}
\omterm{def:m1:pnl-and-positions:exposure}{Exposition brute} $1.0 + 1.2 + 1.5 = \$3.7$ millions, soit $1.85\times$ les fonds propres~; \omterm{def:m1:pnl-and-positions:exposure}{exposition nette} $2.2 - 1.5
= \$0.7$ million, $0.35\times$.
\end{solution}

\begin{solution}{exo:m1:pnl-and-positions:2}
Après le premier achat $\bar p = 20.00$~; après le second $\bar p =
(500\times20.00 + 300\times20.40)/800 = 20.15$. La vente réalise $600\times
0.35 = \$210$ et laisse 200 actions au \omterm{def:m1:pnl-and-positions:realised}{coût moyen} $\bar p = 20.15$. \omterm{def:m1:pnl-and-positions:realised}{P\&L latent} à 20.30~:
$200 \times 0.15 = \$30$.
\end{solution}

\begin{solution}{exo:m1:pnl-and-positions:3}
Trésorerie $-10\,000 - 6\,120 + 12\,300 = -\$3\,820$~; $-3\,820 + 200\times20.30 =
\$240 = 210 + 30$.
\end{solution}

\begin{solution}{exo:m1:pnl-and-positions:4}
Vente à découvert de 1\,000 à 80.00~; vente à découvert de 2\,000 au $\bar p = 80.50$. L'achat de
3\,000 clôture les 2\,000, réalisant $2\,000\times(80.50-79.50) = \$2\,000$,
et ouvre une \omterm{def:m1:pnl-and-positions:position}{position} longue de 1\,000 à 79.50. À 79.80, le \omterm{def:m1:pnl-and-positions:realised}{P\&L latent} vaut 300~\$ et le
total 2\,300~\$.
\end{solution}

\begin{solution}{exo:m1:pnl-and-positions:5}
Prix~: $50\,000\times(-10)\times0.80 = -\pounds400\,000$. Devise~:
$50\,000\times200\times0.04 = +\pounds400\,000$. Terme croisé~:
$50\,000\times(-10)\times0.04 = -\pounds20\,000$. Total $-\pounds20\,000$~:
les deux grands termes s'annulent et le terme croisé est le résultat.
\end{solution}

\begin{solution}{exo:m1:pnl-and-positions:6}
\omterm{def:m1:pnl-and-positions:position}{Position} reportée~: $-20\,000\times(29.70-30.00) = +\$6\,000$. Nouvelles transactions~:
$8\,000\times0.10 + (-5\,000)\times(-0.20) = +\$1\,800$. Frais $-260$~; \omterm{def:m1:financing:seclending}{frais
d'emprunt} $-120$~; intérêts créditeurs sur le produit de vente $+65$. Total $+\$7\,485$. \omterm{def:m1:pnl-and-positions:position}{Position} de clôture
$-17\,000$.
\end{solution}

\begin{solution}{exo:m1:pnl-and-positions:7}
Tenir une file de lots signés~; une exécution de signe opposé consomme des lots depuis la
tête de file, en réalisant $\pm n(p - p_{\text{lot}})$ avec le signe du lot, et
tout reliquat devient un nouveau lot. Pour les exécutions de l'exercice~4~: $1\,000\times0.50 +
1\,000\times1.50 = \$2\,000$, soit le même résultat que le \omterm{def:m1:pnl-and-positions:realised}{coût moyen}. Les deux conventions
s'accordent dès qu'une \omterm{def:m1:pnl-and-positions:position}{position} est clôturée \emph{entièrement}, puisque alors chaque lot est
consommé~; elles ne diffèrent que pour les clôtures partielles.
\end{solution}

\begin{solution}{exo:m1:pnl-and-positions:8}
Le desk clôture ses \omterm{def:m1:pnl-and-positions:position}{positions} gagnantes et garde ouvertes ses \omterm{def:m1:pnl-and-positions:position}{positions} perdantes~: chaque transaction clôturée est
un gain, le \omterm{def:m1:pnl-and-positions:realised}{P\&L réalisé} est positif chaque mois, et les pertes s'accumulent dans
la ligne du \omterm{def:m1:pnl-and-positions:realised}{P\&L latent}, que le rapport omet. La comptabilité par lots, avec un choix libre
du lot à clôturer rend cela le plus aisé, mais toute convention permet
de le faire. La \Cref{prop:m1:pnl-and-positions:identity} y échappe~: le \emph{total}
P\&L valorisé à un prix indépendant, publié chaque jour, fait apparaître la perte dès le
premier mois. Le \omterm{def:m1:pnl-and-positions:realised}{P\&L réalisé} et le taux de réussite décrivent le moment où un trader
a choisi de clôturer, et non le montant gagné.
\end{solution}

\begin{solution}{pb:m1:pnl-and-positions:1}
\textbf{1.} $12\,000 \to 5\,000 \to -4\,000 \to 2\,000 \to 12\,000$.
\textbf{2.} $592\,200 + 764\,100 - 505\,200 - 837\,000 = +\$14\,100$.
\textbf{3.} $32\,000 \times 0.0012 = \$38.40$.
\textbf{4.} Entre la deuxième et la troisième exécution, vente à découvert de 4\,000.
\textbf{5.} La \omterm{def:m1:pnl-and-positions:position}{position} redevient de 12\,000 actions, donc le P\&L brut vaut
$14\,100 + 12\,000(m - 84.00)$. À 83.55~: 8\,700~\$.
\textbf{6.} Cours médian 83.75~: 11\,100~\$.
\textbf{7.} Cours de l'enchère 83.76~: 11\,220~\$.
\textbf{8.} Une \omterm{def:m1:pnl-and-positions:position}{position} longue est clôturée au cours acheteur, 83.70~: 10\,500~\$.
\textbf{9.} Le cours de l'enchère~: il est public, unique et ne peut être influencé
par la dernière petite transaction du trader~; l'écart de 2\,520~\$ entre le plus haut et
le plus bas de ces valorisations reste, lui, à son choix.
\textbf{10.} $12\,000\times(83.76-84.00) = -\$2\,880$.
\textbf{11.} $7\,000\times0.84 = 5\,880$; $9\,000\times1.14 = 10\,260$;
$6\,000\times(-0.44) = -2\,640$; $10\,000\times0.06 = 600$: $+\$14\,100$.
\textbf{12.} $-2\,880 + 14\,100 - 38.40 - 310 = \$10\,871.60$.
\textbf{13.} Brut $11\,220 = -2\,880 + 14\,100$~: rien n'est inexpliqué.
\textbf{14.} Oui, mais regardons de plus près~: presque tout le gain de trading provient de
les deux ventes effectuées au-dessus du cours de clôture, c'est-à-dire du fait d'avoir été \emph{moins
longue} pendant que le cours baissait. Les deux lignes relèvent d'une même décision --- l'idée que le
titre baisserait --- et le second achat en a rendu une partie.
\textbf{15.} Vente de 7\,000~: \omterm{def:m1:pnl-and-positions:realised}{P\&L réalisé} $7\,000\times2.10 = 14\,700$, \omterm{def:m1:pnl-and-positions:position}{position} longue de 5\,000 à
82.50. Vente de 9\,000~: clôture 5\,000 ($+12\,000$, total 26\,700) et ouvre une \omterm{def:m1:pnl-and-positions:position}{position} courte
de 4\,000 à 84.90. Achat de 6\,000~: clôture 4\,000 ($+2\,800$, total 29\,500) et ouvre
une \omterm{def:m1:pnl-and-positions:position}{position} longue de 2\,000 à 84.20. Achat de 10\,000~: \omterm{def:m1:pnl-and-positions:position}{position} longue de 12\,000 au $\bar p = 83.783$.
\textbf{16.} $12\,000\times(83.76-83.783) = -\$280$.
\textbf{17.} $29\,500 - 280 = \$29\,220$ est le P\&L \emph{depuis que la \omterm{def:m1:pnl-and-positions:position}{position}
a été ouverte}. Retrancher le \omterm{def:m1:pnl-and-positions:realised}{P\&L latent} reporté de la veille,
$12\,000\times(84.00-82.50) = \$18\,000$~: 11\,220~\$.
\textbf{18.} Une exécution de 1\,000 actions absente des registres du desk (une
exécution tardive ou saisie à la main, ou une exécution destinée au compte d'un autre desk
imputée ici) --- comparer la liste d'exécution du courtier au journal des exécutions,
par identifiant d'exécution. Viennent ensuite une opération sur titres ou un transfert traité par le
courtier et non par le desk. Enfin, une erreur du courtier.
\textbf{19.} \textbf{10\,872~\$}, net, au cours de clôture officiel~: $+14\,100$ provenant
des transactions du jour, $-2\,880$ sur la \omterm{def:m1:pnl-and-positions:position}{position} reportée, $-348$ de frais et
de financement.
\textbf{20.} Une politique de valorisation écrite --- cours de clôture officiel pour le
P\&L quotidien, net de tous les coûts --- chaque autre chiffre étant étiqueté comme une estimation.
\end{solution}

\begin{solution}{iq:m1:pnl-and-positions:1}
Le \omterm{def:m1:pnl-and-positions:realised}{P\&L réalisé} est le gain sur la part de la \omterm{def:m1:pnl-and-positions:position}{position} qui a été clôturée,
mesuré par rapport à une convention de coût~; le \omterm{def:m1:pnl-and-positions:realised}{P\&L latent} est le gain valorisé sur ce qui
reste ouvert. La convention déplace des montants entre les deux et entre exercices
fiscaux~; elle ne peut modifier le total, qui vaut la trésorerie plus la \omterm{def:m1:pnl-and-positions:position}{position} multipliée par le cours
de valorisation et ne dépend pas de l'ordre dans lequel les achats ont été effectués.
\iqlookfor{l'invariant énoncé sous forme de formule.}
\end{solution}

\begin{solution}{iq:m1:pnl-and-positions:2}
Les nombres à virgule flottante binaire ne peuvent pas représenter exactement la plupart des prix décimaux~; les sommes de
nombreuses exécutions accumulent des erreurs qui dépendent de l'ordre des additions, si bien que deux
systèmes ne s'accordent pas d'un centime et que les rapprochements échouent. Il faut utiliser des entiers exprimés dans une
unité monétaire mineure fixe (en ticks, ou en dix-millièmes de la devise) pour les prix et
la trésorerie, un type décimal dans les couches de reporting, et le calcul à virgule flottante seulement pour les
statistiques dérivées, comme un \omterm{def:m1:pnl-and-positions:realised}{coût moyen}.
\iqlookfor{la dépendance à l'ordre et le rapprochement, et pas seulement l'«~arrondi~».}
\end{solution}

\begin{solution}{iq:m1:pnl-and-positions:3}
Au cours médian, \$500\,000~; pour liquider, au cours acheteur~: \$499\,900. La dernière transaction n'est
pas pertinente~: elle est périmée et éloignée du marché. Pour deux millions d'actions, aucun cours affiché
ne s'applique~: la valeur est le cours médian diminué du coût attendu de la vente
d'une taille pareille, un impact en racine carrée d'au moins quelques dizaines de points de base, et une
valorisation prudente intègre cette réserve.
\iqlookfor{ignorer la dernière transaction, et la valorisation dépendant de la taille.}
\end{solution}

\begin{solution}{iq:m1:pnl-and-positions:4}
Le registre du courtier fait foi juridiquement. Il faut comparer exécution par exécution, par
identifiant, la liste du courtier avec le journal des exécutions de la stratégie et le
drop copy~: une exécution absente de notre côté (accusé de réception perdu, exécution partielle
après une demande d'annulation, exécution pendant une reconnexion) est la cause habituelle. Le
système doit construire sa \omterm{def:m1:pnl-and-positions:position}{position} à partir du drop copy, la rapprocher en continu
de sa propre estimation, et, en cas d'écart, cesser de coter et alerter~: une stratégie
qui opère sur une \omterm{def:m1:pnl-and-positions:position}{position} fausse se couvre dans le mauvais sens.
\iqlookfor{le drop copy comme source de vérité, et l'arrêt du trading comme
réflexe.}
\end{solution}

\begin{solution}{iq:m1:pnl-and-positions:5}
Une décomposition du P\&L en causes --- \omterm{def:m1:pnl-and-positions:position}{positions} portées par facteur de risque,
nouvelles transactions, \omterm{def:m1:pnl-and-positions:carry}{carry}, frais, devise --- dont la somme égale le total.
Un P\&L non expliqué, faible et aléatoire, relève du bruit des approximations~;
s'il est persistant ou croissant, c'est que le modèle de risque ne décrit pas le book (un
facteur de risque manquant, des valorisations erronées ou des paramètres périmés), ou que des transactions
manquent ou sont mal enregistrées. Dans les deux cas, le desk ignore à quoi il est
exposé.
\iqlookfor{le P\&L non expliqué comme signal de contrôle, et non comme une simple gêne comptable.}
\end{solution}

\begin{solution}{iq:m1:pnl-and-positions:6}
Il faut conserver le journal d'événements immuable~: chaque exécution, chaque annulation de transaction et chaque correction comme
événement propre, avec l'identifiant d'exécution, l'horodatage de la bourse, l'horodatage de réception
et une référence à l'événement qu'il corrige. Tout le reste se déduit~: la quantité,
la trésorerie et les frais sont des sommes sur le journal, indépendantes de l'ordre, comme
le montre la \cref{prop:m1:pnl-and-positions:identity}, de sorte que les exécutions tardives ou hors d'ordre
n'exigent aucun traitement particulier pour le total. Les grandeurs dépendantes de l'ordre (\omterm{def:m1:pnl-and-positions:realised}{coût moyen},
\omterm{def:m1:pnl-and-positions:realised}{P\&L réalisé}) sont recalculées en rejouant le journal dans l'ordre de l'heure de la bourse
à partir du dernier point de contrôle précédant l'événement corrigé. Une annulation de transaction est un
événement d'annulation, jamais une suppression, afin que chaque chiffre publié puisse être reproduit
à tout instant.
\iqlookfor{le sourcing par événements, et l'observation que le total exact
ne dépend pas de l'ordre.}
\end{solution}
